\documentclass[10pt,a4paper]{amsart}
\usepackage{amsmath,amssymb,amsthm,graphicx}
% Retained minimal preamble for the Stacks Project chapter flat.tex (More on Flatness).
% Verbatim excerpt of the chapter preamble: Xy-pic setup (native preamble lines 14-18) and the native
% theorem declarations with the mathematical macros (native preamble lines 154-217). Copied mechanically;
% nothing added, removed or redefined. The wrapper class and the amsmath/amssymb/amsthm/graphicx base
% packages are supplied by the assembler, so the class/externaldocument part of the native preamble is omitted.
% For commutative diagrams we use Xy-pic
\usepackage[all]{xy}

% We use 2cell for 2-commutative diagrams.
\xyoption{2cell}
% Theorem environments.
%
\theoremstyle{plain}
\newtheorem{theorem}[subsection]{Theorem}
\newtheorem{proposition}[subsection]{Proposition}
\newtheorem{lemma}[subsection]{Lemma}

\theoremstyle{definition}
\newtheorem{definition}[subsection]{Definition}
\newtheorem{example}[subsection]{Example}
\newtheorem{exercise}[subsection]{Exercise}
\newtheorem{situation}[subsection]{Situation}

\theoremstyle{remark}
\newtheorem{remark}[subsection]{Remark}
\newtheorem{remarks}[subsection]{Remarks}

\numberwithin{equation}{subsection}

% Macros
%
\def\lim{\mathop{\mathrm{lim}}\nolimits}
\def\colim{\mathop{\mathrm{colim}}\nolimits}
\def\Spec{\mathop{\mathrm{Spec}}}
\def\Hom{\mathop{\mathrm{Hom}}\nolimits}
\def\Ext{\mathop{\mathrm{Ext}}\nolimits}
\def\SheafHom{\mathop{\mathcal{H}\!\mathit{om}}\nolimits}
\def\SheafExt{\mathop{\mathcal{E}\!\mathit{xt}}\nolimits}
\def\Sch{\mathit{Sch}}
\def\Mor{\mathop{\mathrm{Mor}}\nolimits}
\def\Ob{\mathop{\mathrm{Ob}}\nolimits}
\def\Sh{\mathop{\mathit{Sh}}\nolimits}
\def\NL{\mathop{N\!L}\nolimits}
\def\CH{\mathop{\mathrm{CH}}\nolimits}
\def\proetale{{pro\text{-}\acute{e}tale}}
\def\etale{{\acute{e}tale}}
\def\QCoh{\mathit{QCoh}}
\def\Ker{\mathop{\mathrm{Ker}}}
\def\Im{\mathop{\mathrm{Im}}}
\def\Coker{\mathop{\mathrm{Coker}}}
\def\Coim{\mathop{\mathrm{Coim}}}

% Boxtimes
%
\DeclareMathSymbol{\boxtimes}{\mathbin}{AMSa}{"02}

%
% Macros for moduli stacks/spaces
%
\def\QCohstack{\mathcal{QC}\!\mathit{oh}}
\def\Cohstack{\mathcal{C}\!\mathit{oh}}
\def\Spacesstack{\mathcal{S}\!\mathit{paces}}
\def\Quotfunctor{\mathrm{Quot}}
\def\Hilbfunctor{\mathrm{Hilb}}
\def\Curvesstack{\mathcal{C}\!\mathit{urves}}
\def\Polarizedstack{\mathcal{P}\!\mathit{olarized}}
\def\Complexesstack{\mathcal{C}\!\mathit{omplexes}}
% \Pic is the operator that assigns to X its picard group, usage \Pic(X)
% \Picardstack_{X/B} denotes the Picard stack of X over B
% \Picardfunctor_{X/B} denotes the Picard functor of X over B
\def\Pic{\mathop{\mathrm{Pic}}\nolimits}
\def\Picardstack{\mathcal{P}\!\mathit{ic}}
\def\Picardfunctor{\mathrm{Pic}}
\def\Deformationcategory{\mathcal{D}\!\mathit{ef}}

% --- appended 2026-10-12 (Chief): the project own definitions of the reference/slogan/history
% display environments, copied mechanically from the shared Stacks preamble lines 8-12
% (\usepackage{verbatim} plus the three \newenvironment lines, which the chapter preamble keeps
% verbatim from the project preamble). Nothing added, removed or redefined: these environments
% carry no mathematical content and the project itself typesets them as comments. Without these
% three lines a counted statement whose own text uses them fails with Environment slogan undefined.
\usepackage{verbatim}
\newenvironment{reference}{\comment}{\endcomment}
\newenvironment{slogan}{\comment}{\endcomment}
\newenvironment{history}{\comment}{\endcomment}
\usepackage{cleveref}
\crefname{originalreference}{original reference}{original references}
\Crefname{originalreference}{Original reference}{Original references}
\makeatletter
\newlabel{situation-existence}{{0CTC}{source}{}{sourcepointer-1}{}}
\newlabel{situation-existence@cref}{{[originalreference][0][]0CTC}{[1][0][]source}{}{}{}}
\makeatother
\crefname{originallemma}{original lemma}{original lemmas}
\Crefname{originallemma}{Original lemma}{Original lemmas}
\makeatletter
\newlabel{lemma-compute-what-it-should-be}{{0CTD}{source}{}{sourcepointer-2}{}}
\newlabel{lemma-compute-what-it-should-be@cref}{{[originallemma][0][]0CTD}{[1][0][]source}{}{}{}}
\makeatother
\crefname{originallemma}{original lemma}{original lemmas}
\Crefname{originallemma}{Original lemma}{Original lemmas}
\makeatletter
\newlabel{lemma-compute-against-perfect}{{0CTE}{source}{}{sourcepointer-3}{}}
\newlabel{lemma-compute-against-perfect@cref}{{[originallemma][0][]0CTE}{[1][0][]source}{}{}{}}
\makeatother
\crefname{originallemma}{original lemma}{original lemmas}
\Crefname{originallemma}{Original lemma}{Original lemmas}
\makeatletter
\newlabel{lemma-relative-pseudo-coherence}{{0CTF}{source}{}{sourcepointer-4}{}}
\newlabel{lemma-relative-pseudo-coherence@cref}{{[originallemma][0][]0CTF}{[1][0][]source}{}{}{}}
\makeatother
\crefname{originallemma}{original lemma}{original lemmas}
\Crefname{originallemma}{Original lemma}{Original lemmas}
\makeatletter
\newlabel{lemma-compute-over-affine}{{0CTG}{source}{}{sourcepointer-5}{}}
\newlabel{lemma-compute-over-affine@cref}{{[originallemma][0][]0CTG}{[1][0][]source}{}{}{}}
\makeatother
\crefname{originallemma}{original lemma}{original lemmas}
\Crefname{originallemma}{Original lemma}{Original lemmas}
\makeatletter
\newlabel{lemma-finitely-presented}{{0CTH}{source}{}{sourcepointer-6}{}}
\newlabel{lemma-finitely-presented@cref}{{[originallemma][0][]0CTH}{[1][0][]source}{}{}{}}
\makeatother
\crefname{originallemma}{original lemma}{original lemmas}
\Crefname{originallemma}{Original lemma}{Original lemmas}
\makeatletter
\newlabel{lemma-proper-support}{{0CTI}{source}{}{sourcepointer-7}{}}
\newlabel{lemma-proper-support@cref}{{[originallemma][0][]0CTI}{[1][0][]source}{}{}{}}
\makeatother
\crefname{originallemma}{original lemma}{original lemmas}
\Crefname{originallemma}{Original lemma}{Original lemmas}
\makeatletter
\newlabel{lemma-proper-pure}{{05K3}{source}{}{sourcepointer-8}{}}
\newlabel{lemma-proper-pure@cref}{{[originallemma][0][]05K3}{[1][0][]source}{}{}{}}
\makeatother
\crefname{originaltheorem}{original theorem}{original theorems}
\Crefname{originaltheorem}{Original theorem}{Original theorems}
\makeatletter
\newlabel{theorem-flat-dimension-n-representable}{{05UG}{source}{}{sourcepointer-9}{}}
\newlabel{theorem-flat-dimension-n-representable@cref}{{[originaltheorem][0][]05UG}{[1][0][]source}{}{}{}}
\makeatother
\crefname{originallemma}{original lemma}{original lemmas}
\Crefname{originallemma}{Original lemma}{Original lemmas}
\makeatletter
\newlabel{lemma-when-universal-flattening}{{05UH}{source}{}{sourcepointer-10}{}}
\newlabel{lemma-when-universal-flattening@cref}{{[originallemma][0][]05UH}{[1][0][]source}{}{}{}}
\makeatother
\crefname{originallemma}{original lemma}{original lemmas}
\Crefname{originallemma}{Original lemma}{Original lemmas}
\makeatletter
\newlabel{lemma-monomorphism-isomorphism}{{0CTJ}{source}{}{sourcepointer-11}{}}
\newlabel{lemma-monomorphism-isomorphism@cref}{{[originallemma][0][]0CTJ}{[1][0][]source}{}{}{}}
\makeatother
\title{Grothendieck existence: a compatible system of modules over the thickenings comes from a finitely presented flat module with proper support}
\author{The Stacks Project Authors}
\date{stacks-project:flat.tex:local-snapshot-20261011; GNU FDL 1.2 or later}
\begin{document}
\maketitle
\noindent\textit{Editorial note.} All delivered mathematical text is the author's own, copied from the pinned Stacks Project chapter file flat.tex; no statement or proof step was rewritten, reordered or removed.\par\medskip
\noindent\textit{Editorial note.} The preamble is the chapter's own retained minimal preamble (sources/stacks-flat/retained-preamble.tex), a verbatim excerpt of the native Xy-pic setup and of the native theorem declarations and mathematical macros, followed by the project own definitions of the reference, slogan and history display environments (verbatim package plus three newenvironment lines) that the retained file carries at its end and that the project itself typesets as comments. No macro was invented or redefined, and the wrapper class and the amsmath/amssymb/amsthm/graphicx base packages are supplied by the assembler.\par\medskip
\noindent\textit{Editorial note.} The counted statement and proof contain no cross reference, no citation and no included file, so the delivered body is self contained.\par\medskip
\noindent\textit{Editorial note.} Canonical location of the counted proposition: Stacks tag 0CTK (sources/stacks-flat/pinned-original-source.tex line 8319).\par\medskip
\noindent\textit{Allowed standard prerequisites (editorial).}\begin{itemize}
\item Construction of the module from the compatible system. Same chapter, Lemmas labelled lemma-compute-what-it-should-be, lemma-compute-against-perfect, lemma-relative-pseudo-coherence, lemma-compute-over-affine, lemma-finitely-presented and lemma-proper-support, cited at proof lines 8328-8334. Produces W and the finitely presented module F on W with proper support and the required reductions.
\item Universal purity of the module relative to the base. Same chapter, Lemma labelled lemma-proper-pure, cited at proof lines 8340-8342. Makes the universal flattening available.
\item Representability of the universal flattening and its monomorphism property. Same chapter, Theorem labelled theorem-flat-dimension-n-representable and Lemma labelled lemma-when-universal-flattening, cited at proof lines 8342-8346. Supplies S' -> Spec(A) as a monomorphism of finite presentation.
\item The monomorphism criterion forcing S' to be the base. Same chapter, Lemma labelled lemma-monomorphism-isomorphism, cited at proof lines 8350-8352. Yields S' = Spec(A) and hence flatness of F over A.
\end{itemize}\medskip
\ifdefined\hypertarget\hypertarget{sourcepointer-1}{}\fi
\noindent\textit{Source reference.} Original reference 0CTC: sources/stacks-flat/pinned-original-source.tex, source lines 8036--8036. Same-chapter pointer into the author environment situation at source line 8036; invoked at proof line None of the counted proof. It is retained only as a navigation pointer; every load-bearing step is inside the delivered span.\par\medskip
\ifdefined\hypertarget\hypertarget{sourcepointer-2}{}\fi
\noindent\textit{Source reference.} Original lemma 0CTD: sources/stacks-flat/pinned-original-source.tex, source lines 8062--8062. Same-chapter pointer into the author environment lemma at source line 8062; invoked at proof line 8329 of the counted proof. It is retained only as a navigation pointer; every load-bearing step is inside the delivered span.\par\medskip
\ifdefined\hypertarget\hypertarget{sourcepointer-3}{}\fi
\noindent\textit{Source reference.} Original lemma 0CTE: sources/stacks-flat/pinned-original-source.tex, source lines 8079--8079. Same-chapter pointer into the author environment lemma at source line 8079; invoked at proof line 8330 of the counted proof. It is retained only as a navigation pointer; every load-bearing step is inside the delivered span.\par\medskip
\ifdefined\hypertarget\hypertarget{sourcepointer-4}{}\fi
\noindent\textit{Source reference.} Original lemma 0CTF: sources/stacks-flat/pinned-original-source.tex, source lines 8128--8128. Same-chapter pointer into the author environment lemma at source line 8128; invoked at proof line 8331 of the counted proof. It is retained only as a navigation pointer; every load-bearing step is inside the delivered span.\par\medskip
\ifdefined\hypertarget\hypertarget{sourcepointer-5}{}\fi
\noindent\textit{Source reference.} Original lemma 0CTG: sources/stacks-flat/pinned-original-source.tex, source lines 8145--8145. Same-chapter pointer into the author environment lemma at source line 8145; invoked at proof line 8332 of the counted proof. It is retained only as a navigation pointer; every load-bearing step is inside the delivered span.\par\medskip
\ifdefined\hypertarget\hypertarget{sourcepointer-6}{}\fi
\noindent\textit{Source reference.} Original lemma 0CTH: sources/stacks-flat/pinned-original-source.tex, source lines 8181--8181. Same-chapter pointer into the author environment lemma at source line 8181; invoked at proof line 8333 of the counted proof. It is retained only as a navigation pointer; every load-bearing step is inside the delivered span.\par\medskip
\ifdefined\hypertarget\hypertarget{sourcepointer-7}{}\fi
\noindent\textit{Source reference.} Original lemma 0CTI: sources/stacks-flat/pinned-original-source.tex, source lines 8245--8245. Same-chapter pointer into the author environment lemma at source line 8245; invoked at proof line 8334 of the counted proof. It is retained only as a navigation pointer; every load-bearing step is inside the delivered span.\par\medskip
\ifdefined\hypertarget\hypertarget{sourcepointer-8}{}\fi
\noindent\textit{Source reference.} Original lemma 05K3: sources/stacks-flat/pinned-original-source.tex, source lines 5016--5016. Same-chapter pointer into the author environment lemma at source line 5016; invoked at proof line 8341 of the counted proof. It is retained only as a navigation pointer; every load-bearing step is inside the delivered span.\par\medskip
\ifdefined\hypertarget\hypertarget{sourcepointer-9}{}\fi
\noindent\textit{Source reference.} Original theorem 05UG: sources/stacks-flat/pinned-original-source.tex, source lines 7895--7895. Same-chapter pointer into the author environment theorem at source line 7895; invoked at proof line 8343 of the counted proof. It is retained only as a navigation pointer; every load-bearing step is inside the delivered span.\par\medskip
\ifdefined\hypertarget\hypertarget{sourcepointer-10}{}\fi
\noindent\textit{Source reference.} Original lemma 05UH: sources/stacks-flat/pinned-original-source.tex, source lines 7989--7989. Same-chapter pointer into the author environment lemma at source line 7989; invoked at proof line 8344 of the counted proof. It is retained only as a navigation pointer; every load-bearing step is inside the delivered span.\par\medskip
\ifdefined\hypertarget\hypertarget{sourcepointer-11}{}\fi
\noindent\textit{Source reference.} Original lemma 0CTJ: sources/stacks-flat/pinned-original-source.tex, source lines 8274--8274. Same-chapter pointer into the author environment lemma at source line 8274; invoked at proof line 8351 of the counted proof. It is retained only as a navigation pointer; every load-bearing step is inside the delivered span.\par\medskip
\begin{theorem}[Grothendieck Existence Theorem]
\label{theorem-existence}
In Situation \ref{situation-existence}
there exists a finitely presented $\mathcal{O}_X$-module
$\mathcal{F}$, flat over $A$, with support proper over $A$,
such that
$\mathcal{F}_n = \mathcal{F} \otimes_{\mathcal{O}_X} \mathcal{O}_{X_n}$
for all $n$ compatibly with the maps $\varphi_n$.
\end{theorem}

\begin{proof}
Apply Lemmas \ref{lemma-compute-what-it-should-be},
\ref{lemma-compute-against-perfect},
\ref{lemma-relative-pseudo-coherence},
\ref{lemma-compute-over-affine},
\ref{lemma-finitely-presented}, and
\ref{lemma-proper-support}
to get an open subscheme $W \subset X$ containing all points
lying over $\Spec(A_n)$
and a finitely presented $\mathcal{O}_W$-module $\mathcal{F}$
whose support is proper over $A$ with
$\mathcal{F}_n = \mathcal{F} \otimes_{\mathcal{O}_W} \mathcal{O}_{X_n}$
for all $n \geq 1$. (This makes sense as $X_n \subset W$.)
By Lemma \ref{lemma-proper-pure} we see that $\mathcal{F}$
is universally pure relative to $\Spec(A)$.
By Theorem \ref{theorem-flat-dimension-n-representable}
(for explanation, see Lemma \ref{lemma-when-universal-flattening})
there exists a universal flattening $S' \to \Spec(A)$
of $\mathcal{F}$ and moreover the morphism $S' \to \Spec(A)$
is a monomorphism of finite presentation.
Since the base change of $\mathcal{F}$ to $\Spec(A_n)$
is $\mathcal{F}_n$ we find that $\Spec(A_n) \to \Spec(A)$
factors (uniquely) through $S'$ for each $n$.
By Lemma \ref{lemma-monomorphism-isomorphism}
we see that $S' = \Spec(A)$.
This means that $\mathcal{F}$ is flat over $A$.
Finally, since the scheme theoretic support $Z$ of $\mathcal{F}$
is proper over $\Spec(A)$, the morphism $Z \to X$ is closed.
Hence the pushforward $(W \to X)_*\mathcal{F}$ is supported
on $W$ and has all the desired properties.
\end{proof}

\end{document}
